\documentclass[11pt,a4paper]{ltjsarticle}
\usepackage[margin=24mm]{geometry}
\usepackage[haranoaji]{luatexja-preset}
\usepackage{amsmath,amssymb,amsthm,mathtools}
\usepackage{booktabs,longtable,array}
\usepackage{xcolor,hyperref}
\hypersetup{colorlinks=true,linkcolor=blue!45!black,urlcolor=blue!45!black,
pdftitle={三重超平面と接二次超曲面を退化因子にもつPoisson構造},
pdfauthor={Poisson P4 research project}}
\setlength{\emergencystretch}{3em}
\newcommand{\CC}{\mathbf C}
\newcommand{\QQ}{\mathbf Q}
\newcommand{\PP}{\mathbf P}
\renewcommand{\AA}{\mathbf A}
\newcommand{\GG}{\mathbf G}
\newcommand{\cP}{\mathcal P_4}
\newcommand{\cC}{\mathcal C_{\mathrm{tq}}}
\newcommand{\Sym}{\operatorname{Sym}}
\DeclareMathOperator{\GL}{GL}
\DeclareMathOperator{\PGL}{PGL}
\DeclareMathOperator{\Sp}{Sp}
\DeclareMathOperator{\rank}{rank}
\DeclareMathOperator{\diver}{div}
\DeclareMathOperator{\Spec}{Spec}
\DeclareMathOperator{\Stab}{Stab}
\newtheorem{theorem}{定理}[section]
\newtheorem{proposition}[theorem]{命題}
\newtheorem{lemma}[theorem]{補題}
\newtheorem{corollary}[theorem]{系}
\theoremstyle{remark}
\newtheorem{remark}[theorem]{注意}
\renewcommand{\proofname}{証明}
\title{三重超平面と接二次超曲面を退化因子にもつ\\
射影四空間の Poisson 構造\\[2mm]
\large 明示的な25次元成分とその一次ジェット像}
\author{Poisson P4 研究プロジェクト・計算証明付き研究草稿}
\date{2026年11月3日・v63}

\begin{document}
\maketitle
\begin{abstract}
射影四空間の Poisson スキームに、明示的な三変数族の射影線形群による
飽和閉包として25次元の既約成分を構成する。
一般の元の実際の退化五次因子は $3H+Q$ であり、$Q$ は滑らかな二次超曲面、
$H$ はその接超平面である。
モジュラー流の線形作用素は一般に Jordan 型 $(3,1,1)$ をもち、
連結射影安定化群は加法群である。
追加の射影せん断を含む有限座標変換により、一般の元を二変数の族に移せる。
この表示は一般軌道の有限被覆を与えるが、残る有限同値関係の完全分類は行わない。

成分であることの証明には、パラメータ原点での\emph{全} Poisson 方程式の形式横断環を用いる。
126個の正準係数から軌道方向と95個の陰関数を除去し、八変数の全残差形式級数を保つ。
四次までの真の初期式に含まれる21個の斉次式が次元3のイデアルを生成することから、
全射影局所次元の上界25を得る。
一方、元の三変数族と群作用の微分には26次の非零有理小行列式があり、
像の射影次元は25以上である。両者を合わせると、飽和閉包そのものが既約成分となる。

この成分の primitive 一次ジェット像は16次元空間内の15次元既約超曲面であり、
その像上の一般 $S_{61}$ ファイバーの次元は1である。
従来の57成分との相異性を示し、成分数の下界を58に更新する。
さらに全 Jacobi 行列の関数体階数98と二法線の二次障害の階数1から、
一般横断 Artin 環の長さを5と決定する。
点 $(a,b,c)=(1,1,1)$ での全射影完備局所環は
$\CC[[r_1,\ldots,r_{25},u,v]]/(uv,u^3,v^3)$ であり、
Cohen--Macaulay・type 2 である。
一方、パラメータ原点での全局所環は Cohen--Macaulay でない。
全成分の尽くし、一般一次ジェット上の五点問題、被約成分の大域的正規性は結論しない。
\end{abstract}

\section{対象と主定理}

基礎体を $\CC$ とし、$V=\CC^5$、$\PP^4=\PP(V)$ とする。
本稿の対象は
\[
 \cP=\bigl\{[\pi]\in
 \PP H^0(\PP^4,\wedge^2T_{\PP^4}):[\pi,\pi]=0\bigr\}
\]
という射影 Jacobi スキームである。
射影化は非零スカラー倍を除く操作であり、$\PGL_5$ による商は取らない。
「成分」は特記しなければ被約台の既約成分を意味する。
局所次元の上界では被約化前の全方程式を使う。

斉次座標を $(h,x,y,z,w)$ とし、$h=1$ 上で
\begin{equation}\label{eq:family}
\begin{aligned}
 \{x,y\}&=\frac b{36}-\frac{256}3w,\\
 \{x,z\}&=\frac{bc}{48}y+\frac b{48}x-64cwy-64wx,\\
 \{x,w\}&=\frac89,\\
 \{y,z\}&=\frac{ac}{2}y^2+axy-\frac{4a}{3}z
 -\frac{b^2}{1536}w+bw^2-\frac b{64}x^2+64wy-\frac29,\\
 \{y,w\}&=\frac43x,\\
 \{z,w\}&=cxy+x^2+\frac23y
\end{aligned}
\end{equation}
と定める。この反対称二ベクトルを $P_{a,b,c}$ と書く。
右辺は全て二次以下なので斉次化により $\PP^4$ 上の二ベクトルを定める。
後で示すとおり、全ての $(a,b,c)$ で Poisson である。
原点の元を $P_+=P_{0,0,0}$ とし、
\[
 \cC=\overline{\PGL_5\cdot
       \{[P_{a,b,c}]:(a,b,c)\in\AA^3\}}
\]
とおく。閉包は $\cP$ 内で取る。

\begin{theorem}\label{thm:main}
次が成り立つ。
\begin{enumerate}
\item $\cC$ は $\cP$ の25次元既約成分である。
\item $\cC$ の一般の元の実退化因子は $3H+Q$ であり、
 $Q$ は滑らかな二次超曲面、$H$ は $Q$ の接超平面である。
 一般射影軌道の次元は23、連結射影安定化群は $\GG_a$ である。
\item 一般の元は $P_{u,1,v}$ と射影同値である。
 この二変数族は一般軌道空間の有限被覆を与える。
\item 値と perspective を正規化した primitive 一次ジェットの空間
 $W\simeq\AA^{16}$ における像閉包は15次元の既約超曲面である。
 対応する16次元の $S_{61}$ 成分からこの像への一般ファイバーの次元は1である。
\item 全 Jacobi スキームは $\cC$ に沿って一般に非被約である。
 一般横断 Artin 局所環の埋め込み次元は2、長さは5である。
 点 $[P_{1,1,1}]$ での全射影完備局所環は
 \[
 \CC[[r_1,\ldots,r_{25},u,v]]/(uv,u^3,v^3)
 \]
 と同型であり、被約環は滑らか、全環は Cohen--Macaulay・type 2 である。
\item 点 $[P_+]$ での全 Poisson 局所環は Cohen--Macaulay でない。
\end{enumerate}
\end{theorem}

既知の成分との比較には、MPS の40個の成分の識別
\cite{mps} と、研究リポジトリに保存された11個の追加一般階数四成分・六個の
一般階数二成分の構成と相異性を入力とする。
\begin{corollary}\label{cor:58}
上記の既知57成分と $\cC$ は全て相異なる。
従って $\cP$ の既約成分数は58以上、一般階数四成分数は52以上である。
\end{corollary}
これは保存された既知カタログに対する下界である。
上界、分類の完了、文献全体での優先性は主張しない。

\section{二次二ベクトルの射影模型}

\subsection{Euler 商と正準係数}

$X=(X_0,\ldots,X_4)=(h,x,y,z,w)$、
$E=\sum_iX_i\partial_i$ とする。
斉次二次係数をもつ反対称二ベクトルは
\[
 B=\sum_{0\le i<j\le4}B^{ij}(X)\partial_i\wedge\partial_j
 \in \wedge^2 V\otimes\Sym^2V^*
\]
と書け、その係数空間の次元は $10\cdot15=150$ である。
射影化で零となるものは $E\wedge A$（$A$ は線形ベクトル場）であり、
この部分空間の次元は24である。

本稿では
\[
 (\delta B)^j=\sum_i\partial_i B^{ij}
\]
という発散を用い、$\delta B=0$ の正準代表を取る。
斉次二次反対称行列では $\delta B$ は跡零の線形ベクトル場である。
$\delta(E\wedge A)=5A$ が跡零の $A$ について成立するので、
\begin{equation}\label{eq:canonical}
 B^{\mathrm{can}}=B-\frac15E\wedge\delta B
\end{equation}
は元と同じ射影二ベクトルを表す。
従って正準係数空間は126次元である。

\begin{lemma}\label{lem:strict}
正準代表 $B$ について、射影 Poisson 条件は厳密な多項式等式
$[B,B]=0$ と同値である。
その係数は350個の二次方程式として保存できる。
\end{lemma}
\begin{proof}
Schouten 括弧の係数は斉次三次で、
$\binom53\binom73=350$ 個ある。
射影的に零なら $[B,B]=E\wedge C$ と書ける。
発散作用素の BV 恒等式と $\delta B=0$ から $\delta[B,B]=0$ である。
一方、斉次次数3の三ベクトルに対する Euler--発散の恒等式により、
$\delta(E\wedge C)=5C-E\wedge\delta C$ となる。
従って $5E\wedge C=E\wedge E\wedge\delta C=0$ であり、$[B,B]=0$ を得る。
逆向きは直ちに従う。
\end{proof}

\subsection{全パラメータの Jacobi 恒等式}

式\eqref{eq:family}の係数行列を $P=(P^{ij})$ とする。
直接の展開で、四つの相異なる三添字 $i<j<k$ に対して
\begin{equation}\label{eq:affine-jacobi}
 \sum_{\ell=1}^4
 \left(P^{i\ell}\partial_\ell P^{jk}
      +P^{j\ell}\partial_\ell P^{ki}
      +P^{k\ell}\partial_\ell P^{ij}\right)=0
 \quad\text{in }\QQ[a,b,c,x,y,z,w]
\end{equation}
が成立する。
この等式は $N^2=0$ などの追加条件を課していない。
特に $c$ は自由変数のままである。
各括弧を二次斉次化し、$B^{0i}=0$ として得る raw 代表を
式\eqref{eq:canonical}で正準化した全150係数についても、
補題\ref{lem:strict}の全350係数が恒等的に零となる。
検証用データでは $a,b,c$ を $n0,n1,n2$ と表記する。

\section{実退化因子とモジュラー流}

\subsection{三重超平面と滑らかな二次超曲面}

一般階数四二ベクトルの二乗は $\wedge^4T_{\PP^4}\simeq\mathcal O(5)$ の切断であり、
その零因子を\emph{実退化因子}と呼ぶ。
被約台だけを取る操作は行わない。
本族のアフィン Pfaffian を
\[
 F=P^{xy}P^{zw}-P^{xz}P^{yw}+P^{xw}P^{yz}
\]
とすると、三次・四次項は全て打ち消し合い、$F$ は二次になる。
その二次斉次化を $Q$ と書くと
\begin{equation}\label{eq:quadric}
\begin{aligned}
 5184Q={}&2304ac\,y^2+4608a\,xy-6144a\,hz
 -3b^2hw+4608b\,w^2\\
 &{}-72b\,x^2+96b\,hy-1024h^2.
\end{aligned}
\end{equation}
raw 代表の射影二乗は正確に $h^3Q$ となる。
従って実因子は $3H+V(Q)$（$H=V(h)$）である。

\begin{proposition}\label{prop:quadric}
$ab(32a+bc)\ne0$ なら、$Q$ は滑らかで $H$ は $Q$ に接する。
接点は $[0:0:0:1:0]$ である。
\end{proposition}
\begin{proof}
$Q=X^TM_QX$ と書くと
\begin{equation}\label{eq:qdet}
 \det M_Q=\frac{1024}{531441}a^3b(32a+bc).
\end{equation}
従って仮定の下で $Q$ は滑らかである。
$h=0$ への制限は $z$ に依存せず、$(x,y,w)$ に関する小行列式は
\[
 -\frac4{729}ab(32a+bc)
\]
である。よって $Q\cap H$ は階数3の二次錐で、頂点は上記の点である。
そこでの $Q$ の勾配は $(-32a/27,0,0,0,0)$ なので接超平面は $h=0$ となる。
\end{proof}

原点では $Q=-16h^2/81$、従って実因子は $5H$ である。
これは一般の $3H+Q$ の退化である。
一般の二次曲面は階数5なので、階数4以下の二次曲面を伴う従来の三重超平面型とは異なる。

\subsection{指数3の加法群作用}

発散のアフィン表示は
\begin{equation}\label{eq:modular}
 V_{\mathrm{mod}}=\frac{4a}{3}\partial_y+
             \left(ax+acy+\frac b{48}\right)\partial_z.
\end{equation}
これを斉次線形ベクトル場と見た行列を $N$ とすると、非零成分は
\begin{equation}\label{eq:N}
 N_{y h}=\frac{4a}3,\qquad N_{z h}=\frac b{48},\qquad
 N_{z x}=a,\qquad N_{z y}=ac
\end{equation}
だけである。
$N^3=0$ で、$N^2$ の唯一の非零候補は
$(N^2)_{z h}=4a^2c/3$ である。
$ac\ne0$ なら階数は $(\rank N,\rank N^2)=(2,1)$、
Jordan 型は $(3,1,1)$ となる。

\begin{proposition}\label{prop:Ga}
全パラメータで $\exp(tN)=I+tN+t^2N^2/2$ は raw 代表を厳密に保つ。
また
\[
 N^TM_Q+M_QN=0,\qquad
 P\,\nabla F=-F\,V_{\mathrm{mod}}
\]
が成立する。
一般の連結射影安定化群はこのモジュラー加法群に等しい。
\end{proposition}
\begin{proof}
最初の二つは式\eqref{eq:family}--\eqref{eq:N}への代入による有理多項式恒等式である。
最後の恒等式は非退化部分で $V_{\mathrm{mod}}=-P\nabla\log F$ とも読める。
軌道次元については次節の全係数小行列式で一般射影軌道次元23を示す。
$\dim\PGL_5=24$ なので一般射影安定化群の次元は1である。
既に非自明な閉部分群 $\GG_a$ を含むため、その連結成分がこれに等しい。
\end{proof}

\section{軌道の次元と二変数表示}

\subsection{全係数微分の非零小行列式}

$\rho(P)$ を正準二ベクトルの150係数列とする。
行は $(i,j)$ を辞書式順序で並べ、各組について二次単項式 $X_rX_s$（$r\le s$）
を辞書式順序で並べる。
線形行列 $M$ の作用の微分は
\[
 \mathcal L_M B=MB+BM^T-\sum_i(MX)_i\partial_iB
\]
で与えられる。
点 $(a,b,c)=(1,2,3)$ で、25個の $M=E_{ij}$、スカラー方向 $P$、
三つの族方向からなる $150\times29$ 行列を作る。
\begin{lemma}\label{lem:rank}
この行列の階数は26である。
$\GL_5$ 方向のみの階数は23、スカラー方向を加えた階数は24である。
\end{lemma}
\begin{proof}[計算証明]
全係数を $\QQ$ 上で作り、小行列式と全列の従属関係を検証する。
26次小行列式の一つは
\[
 -\frac{76813395244431179776}{1946195068359375}\ne0
\]
となる。その行・列添字を付録\ref{app:rank}に記す。
23次・24次の証人はそれぞれ
\[
 \frac{4800837202776948736}{216243896484375},\qquad
 -\frac{38406697622215589888}{648731689453125}
\]
である。保存証明書には選択列に対する全列の係数行列もあり、
選択列との積が元の全行列になることを検証する。
\end{proof}

非零小行列式は標数零の像の次元下界を与える。
従って $\GL_5\times\GG_m\times\AA^3$ の像錐の次元は26以上、
その射影像閉包 $\cC$ の次元は25以上である。
また、モジュラー加法群により射影軌道次元は一般に23以下であり、
補題\ref{lem:rank}の下界と合わせて一般に23となる。

\subsection{定数項を保つ射影せん断}

$a\ne0$、$r\ne0$、$\epsilon\in\{1,-1\}$ とする。
有限射影変換
\begin{equation}\label{eq:normalization}
\begin{aligned}
 h'&=h,& x'&=rx,& y'&=r^2y,& w'&=\epsilon r w,\\
 z'&=r^3z+\frac{r^3-\epsilon}{6a}h
\end{aligned}
\end{equation}
を $g$ と書く。
全二ベクトルの恒等式として
\begin{equation}\label{eq:parameter-action}
 \epsilon r^{-2}g_*P_{a,b,c}=P_{\epsilon a,\,\epsilon rb,\,c/r}
\end{equation}
が成り立つ。
特に $ab\ne0$ で $\epsilon=1,r=1/b$ と取れば
$P_{a,b,c}$ は $P_{a,1,bc}$ と射影同値になる。
二変数を $(u,v)=(a,bc)$ とおくと、さらに $(u,v)\sim(-u,-v)$ という明示同値がある。

\begin{remark}
式\eqref{eq:normalization}の $h$ 方向のせん断は不可欠である。
対角変換だけでは $\{y,z\}$ の定数 $-2/9$ が変化する。
せん断の定数を $d=(r^3-\epsilon)/(6a)$ とすると、変換後の定数は
\[
 \frac{4\epsilon a}3d-\frac{2\epsilon r^3}9=-\frac29
\]
となる。式\eqref{eq:parameter-action}は六つのアフィン括弧と
斉次化後の全150係数の両方で検証されている。
\end{remark}

一般軌道の次元23と $\dim\cC=25$ が成立すると、Rosenlicht 商の次元は2となる。
上の二変数族はその商を優越し、一般ファイバーは零次元なので有限被覆を与える。
これは他の有限同値の不存在を意味しない。
二変数から完全な軌道不変量を得る問題は残る。

\section{全 Poisson 局所環の次元上界}\label{sec:upper}

前節の下界だけでは、25次元の部分族がより大きな成分に含まれる可能性を除けない。
本節では $P_+$ における\emph{周囲の全} Poisson 局所環に上界25を与える。
特殊な三変数族の方程式だけを調べる操作ではない。

\subsection{八変数の全形式横断模型}

補題\ref{lem:strict}の全350式を $\mathcal J$ と書く。
126次元正準係数空間で $P_+$ を原点に平行移動する。
$\GL_5\times\GG_m$ の軌道微分はこの点では階数23である。
23本の線形切断を選ぶ軌道小行列式は
\[
 -\frac{332306998946228968225951765070086144}{14348907}\ne0
\]
となる。群の23次元局所切断と103次元線形横断面との作用写像はエタールなので、
完備局所環は23個の滑らか変数と横断環の形式冪級数拡大になる。
射影化すると滑らか因子は22個である。

全 $350\times103$ Jacobi 線形化の階数は95である。
95個の従属座標と95本の\emph{元の} Jacobi 方程式を選び、
可逆な95次行列を $B$ と書く。
形式陰関数定理により、残る八変数 $z=(z_0,\ldots,z_7)$ に関する
従属座標の一意な形式級数が存在する。
その代入後も元の残り255本の残差級数を全て保ち、生成するイデアルを $\mathfrak I$ とする。
従って
\begin{equation}\label{eq:full-formal}
 \widehat{\mathcal O}_{\cP,[P_+]}
 \simeq\CC[[t_1,\ldots,t_{22},z_0,\ldots,z_7]]/\mathfrak I.
\end{equation}
これは全次数の形式的同型である。
二次・三次までの打切りを完備局所環と同一視していない。

\subsection{真の初期式を四次まで作る}

全 Jacobi 式は二次なので、陰関数の各斉次項は有限演算で再帰的に得られる。
線形解を $R_1=Tz$ と書き、$R_d$ は選択95式の次数 $d$ 残差に
$-B^{-1}$ を掛けたものとして求める。
この手順で $R_2,R_3,R_4$ を作り、\emph{全350式}に戻す。
選択95式の残差は四次まで零となる。

二次残差の張る空間の階数は5であり、定数倍を調整した基底は
\begin{equation}\label{eq:quadrics}
\begin{aligned}
 A_1&=z_1^2+120z_5z_7-14z_7^2,\\
 15A_2&=15z_1z_3-5z_1z_6+12z_3z_7+8z_4z_5-2z_4z_7
                         +16z_5z_6-8z_6z_7,\\
 3A_3&=3z_1z_4+2z_1z_6+36z_3z_7-16z_4z_5+12z_4z_7
                         -32z_5z_6+8z_6z_7,\\
 4A_4&=4z_1z_5+36z_5z_7-z_7^2,\\
 A_5&=z_7(z_1-20z_5+5z_7)
\end{aligned}
\end{equation}
となる。二次式だけのイデアルは次元5であり、目標の上界には足りない。

二次項の独立な五本の\emph{残差級数}を選び、他の345本からその定数線形結合を引く。
得られる級数 $G_j$ は三次から始まる。
次数3の係数空間で $z_i A_\nu$ と $G_{j,3}$ を合わせた385本の行列を作ると、
階数48、全関係空間の次元337となる。
この\emph{全}関係を級数の組合せに持ち上げて四次項を取り出す。
次数4では $\Sym^2(z)A_\nu$、$z_iG_{j,3}$ とこの337本の繰上がりを合わせる。
その階数は204である。

\begin{lemma}[有限有理証明書]\label{lem:initial}
上記の操作により、5個の二次式、8個の三次式、8個の四次式からなる
$\QQ[z_0,\ldots,z_7]$ の斉次イデアル $L$ が得られる。
全ての生成元は $\mathfrak I$ の真の初期イデアルに属する。
その Hilbert 級数は
\begin{equation}\label{eq:hilbert}
 \operatorname{Hilb}_{\QQ[z]/L}(T)=
 \frac{1+5T+10T^2+2T^3-5T^4-7T^5-7T^6-T^7+4T^8+4T^9}
 {(1-T)^3}.
\end{equation}
特に $\dim\QQ[z]/L=3$、次数は6である。
\end{lemma}
\begin{proof}[計算証明]
選択行、全関係337本、各関係の四次繰上がり、21個の生成元を有理係数で保存する。
元の350式への代入によって、各生成元が真の初期式であることを検証する。
八変数の21式に対する $\QQ$ 上の Gröbner 計算で式\eqref{eq:hilbert}を得る。
分子の $T=1$ での値は6なので約分後の極の位数は3である。
全計算の入力と独立検証の対応は付録\ref{app:cert}に記す。
\end{proof}

\begin{remark}
式\eqref{eq:hilbert}は $L$ の級数であり、未知の全初期イデアルや
全局所環の Hilbert 級数であるとは主張しない。
本稿の上界には $L\subset\operatorname{in}(\mathfrak I)$ だけで十分である。
\end{remark}

\subsection{成分の認識}

\begin{proof}[定理\ref{thm:main}(1)の証明]
補題\ref{lem:initial}と完備局所環の随伴次数環の次元不変性から
\[
 \dim\CC[[z]]/\mathfrak I
 =\dim\CC[z]/\operatorname{in}(\mathfrak I)\le3.
\]
式\eqref{eq:full-formal}により $\dim_{[P_+]}\cP\le22+3=25$ となる。
一方、補題\ref{lem:rank}により $\dim\cC\ge25$ である。

像の始域 $\GL_5\times\GG_m\times\AA^3$ は既約なので $\cC$ は既約である。
$\cC$ を含む $\cP$ の既約成分 $C$ を取ると、$P_+\in\cC\subset C$ である。
有限型既約多様体の閉点における局所次元は全体の次元に等しいから
\[
 25\le\dim\cC\le\dim C\le\dim_{[P_+]}\cP\le25.
\]
全て等号であり、閉既約部分集合 $\cC\subset C$ が同次元なので $\cC=C$ である。
これで成分であることと、全射影局所次元25・実横断環次元3も従う。
\end{proof}

命題\ref{prop:quadric}、\ref{prop:Ga}、式\eqref{eq:parameter-action}から
定理\ref{thm:main}(2),(3)も得られる。
この証明は $P_+$ に他の枝がないことを必要としない。
成分に沿うスキーム的厚みは第\ref{sec:nonreduced}節で別に調べる。

\section{既知の成分との相異性}

五次式の射影空間の中で
\[
 \Sigma_{3,2}=\{[L^3Q]:0\ne L\in V^*,\ 0\ne Q\in\Sym^2V^*\}
\]
を考える。
これは射影多様体 $\PP(V^*)\times\PP(\Sym^2V^*)$ からの固有像なので閉である。
一般階数四の開部分では実五次式への写像は正則である。
式\eqref{eq:quadric}と閉性により、$\cC$ の\emph{全ての階数四点}の実五次式は
$\Sigma_{3,2}$ に入る。

MPS の40成分はそれぞれ、実退化因子が五つの\emph{相異なる}座標超平面
\[
 H_0+H_1+H_2+H_3+H_4
\]
である toric な元を含む\cite{mps}。
この因子は一つの超平面の五重化 $5H$ ではなく、三重線形因子をもたない。
従って $\Sigma_{3,2}$ の外にある。
これらの40成分のどれも $\cC$ に等しくない。

この toric 元が保存図式の実パラメータに属することも、有理係数で直接確認できる。
各図式の biresidue の $4\times4$ 小行列を $B$ とすると、全40件で $\operatorname{Pf}(B)\ne0$ である。
$A=-\operatorname{Pf}(B)B^{-1}$ は実 toric 括弧を与え、
\[
 D=\operatorname{Pf}(A)X_0X_1X_2X_3X_4
   =\operatorname{Pf}(B)X_0X_1X_2X_3X_4
\]
となる。各 $B$ の period 部分空間への帰属も保存基底で検証する。
図式から成分への識別には、MPS の局所 Torelli 定理と射影空間での成分の相異性を用いる。

従来の11個の追加一般階数四成分の射影次元は22,23,24,28のいずれかであり、25のものはない。
従ってそれらとも異なる。
また、一般階数二の六成分は $\pi\wedge\pi=0$ という閉条件に全体が含まれる。
$P_+$ の実五次式は $-16h^5/81\ne0$ なので、本成分はそのいずれとも異なる。
これで系\ref{cor:58}が従う。

ここで「58以上」という結論は、既知57成分の相互相異性を入力としている。
本稿だけで旧57成分の構成全体を再証明するものではない。
新成分の実因子と次元に基づく区別は、その入力の後に独立に検証できる。

\section{一次ジェット像は15次元である}

\subsection{正規化と \texorpdfstring{$S_{61}$}{S61}}

固定点で二ベクトルの値が非退化である開部分を取る。
局所座標で値を固定した標準交代行列 $J$ にし、
原点を固定する四次元の perspective 群で一次項の四つの跡方向を除く。
残る primitive 一次項の空間を $W\simeq\AA^{16}$ とする。
この正規化では残る $\Sp_4$ の作用は商に取らない。

保存された元の三角消去により、Poisson 方程式は
重み $1^{16},2^{45}$ の61座標、134方程式の模型 $S_{61}$ と同値になる。
一般階数四の射影成分 $C$ と対応する $S_{61}$ 成分 $C_{61}$ の間には
\begin{equation}\label{eq:dimension-nine}
 \dim C=9+\dim C_{61}
\end{equation}
が成立する。
9は射影化した非退化値の自由度5と perspective の自由度4の和である。
本節ではこの既存の正規化・イデアル同値を入力とする。
従って $\cC$ に対応する成分の次元は16である。

\subsection{全関数体上の階数計算}

一般評価点 $p=(p_0,p_1,p_2,p_3)$ と局所変数 $t$ に対し
\[
 P(p+t)=C+L(t)+\text{二次項}
\]
と書く。値 $C$ の六係数と一次項 $L$ の24係数を合わせた30係数を取る。
始域の接方向は
\[
 \underbrace{(a,b,c)}_{3},\quad
 \underbrace{p}_{4},\quad
 \underbrace{\mathfrak{gl}_4}_{16},\quad
 \underbrace{\text{perspective}}_{4},\quad
 \underbrace{\text{scalar}}_{1}
\]
の計28方向である。
perspective の列は二次ベクトル場
$V_i=-t_i\sum_jt_j\partial_{t_j}$ の作用を微分して求める。
この全 $30\times28$ 行列を $M_{\mathrm{jet}}$ とする。

\begin{lemma}\label{lem:jet}
体 $K=\QQ(a,b,c,p_0,p_1,p_2,p_3)$ 上で
$\rank M_{\mathrm{jet}}=25$ である。
\end{lemma}
\begin{proof}[計算証明]
まず非零の\emph{有理定数}をピボットとする20回の行列消去を、
七変数の多項式環全体で行う。
残る $10\times8$ 行列に対し $K$ 上で厳密な有理関数消去を行うと階数5となる。
残りの行が全て零であることを関数体で検証する。
独立検証では元の全ジェット行列を別に構成し、SymPy の有理関数体と
FLINT に基づく別の消去実装の両方でこの階数を確認する。
従って標本階数から一般上界を外挿した計算ではない。
下界の別の25次有理小行列式は
\[
 \frac{815115503757040812032}{68630377364883}\ne0
\]
である。
\end{proof}

\begin{proof}[定理\ref{thm:main}(4)の証明]
非退化値 $C$ を動かす $\GL_4$ の方向は独立な六方向を与える。
perspective は値を保ちながら一次項に独立な四方向を与える。
これらの方向を既に全て始域に含めているので、局所的な正規化は
raw ジェット像を六値方向・四跡方向・primitive 像に分ける。
従って primitive 像の次元は正確に
\[
 25-6-4=15
\]
である。
この差し引きで $\Sp_4$ や残るスカラー方向をさらに除いてはいない。

始域が既約なので像閉包も既約であり、$W$ の中の余次元1の閉集合、すなわち既約超曲面である。
式\eqref{eq:dimension-nine}による16次元成分からこの像への一般ファイバーは
$16-15=1$ 次元となる。
\end{proof}

\begin{remark}
この成分は $W$ を優越しない。
従って、十分一般の $\beta\in W$ 上で既知の五つの拡張が全ファイバーを尽くすかという
別の問題への反例にはならない。
像超曲面の展開方程式、一般ファイバーの被約性・既約性、特殊ファイバーの全構造は未決である。
\end{remark}

\section{成分に沿う一般非被約性}\label{sec:nonreduced}

全126正準係数に関する350本の Jacobi 方程式の微分を、
三変数族 $P_{a,b,c}$ に代入した行列を $M_{\mathrm{Jac}}$ とする。
\begin{lemma}\label{lem:generic-jacobian}
体 $\QQ(a,b,c)$ 上で $\rank M_{\mathrm{Jac}}=98$ である。
\end{lemma}
\begin{proof}[計算証明]
全 $350\times126$ 係数を三変数多項式環で構成し、
70個の非零有理定数をピボットとする消去を行う。
恒等的な零行73本・零列10本を除くと $207\times46$ 行列が残る。
この残余を三変数有理関数体上で消去した厳密な階数は28であり、
最終の残余係数が全て零になることを検証する。
従って全体の階数は $70+28=98$ である。
別の独立検証では、元の全350物理 Jacobi 式の係数辞書から微分を作り直し、
全多項式行列、70単元消去、関数体計算への全入力を照合する。
98次非零有理小行列式も別に確認する。
\end{proof}

\begin{proof}[一般非被約性と埋め込み次元]
本族の群・スカラー飽和は $\cC$ に稠密であり、Jacobi 行列の階数は群作用で不変である。
従って補題\ref{lem:generic-jacobian}の階数は成分の一般点に適用できる。
全スキームの一般アフィン接空間次元は $126-98=28$、射影接空間次元は27となる。
一方、成分自体の射影次元は25である。
相異なる他成分との交叉は一般点を含まず、標数零の被約成分は稠密開で滑らかである。
従って全スキームがこの成分に沿って一般被約なら一般接空間次元は25となり、矛盾する。

一般点で局所化した横断零次元環では、接空間の超過分 $27-25=2$ が
極大イデアルをその二乗で割った空間の次元である。
従って埋め込み次元は2、長さは $1+2=3$ 以上である。
\end{proof}

\begin{remark}
この結論は全スキームの非被約性であり、被約成分自体が一般に特異であるという意味ではない。
一点での閉じた Artin 切断の長さから一般長さや平坦性を推定してもいない。
次の補題と有限階数の比較によって、その追加の根拠を与える。
\end{remark}

\subsection{二次障害と閉ファイバーの長さ}

正しいせん断正規化 $b=1$ の稠密開上で、
全 Jacobi 行列の kernel に入る二つの次数1多項式列 $T_1,T_2$ を取れる。
これらは26個の実際の群・スカラー・族方向と合わせて28次元の全 kernel を張る。
元の二次 Jacobi 式に $sT_1+tT_2$ を入れ、$s^2,st,t^2$ の係数列を
$H$ と書く。

\begin{lemma}\label{lem:augmented}
$\QQ(a,c)$ 上で $\rank[J\mid H]=99$ である。
従って実26方向を割った二法線の二次障害の階数は $99-98=1$ となる。
\end{lemma}
\begin{proof}[計算と降下]
元の全350物理 Jacobi 式から $J$ と全三列 $H$ を構成する。
79個の非零有理定数をピボットに消去し、恒等零98行・8列だけを除くと
$173\times42$ 行列が残る。
全二変数有理関数体での残余の階数は20であり、全体は $79+20=99$ となる。
独立監査では物理二次係数辞書から全係数を再構成し、全79操作、
実際の関数体実行入力の全7266分子・分母、最終零行を検査する実装を照合する。
99次の非零有理小行列式も別に検査する。

普遍恒等式 $JT_i=0$ を族方向で微分すると、族方向と $T_i$ との混合 Hessian は
$J$ の像に入る。群・スカラー方向でも同変性から同じことが成り立つ。
実族方向同士の二次項は、全次数 Poisson 族の二階微分によって積分される。
従って $H$ の $J$ の余核への像は、26方向を割った固有の二次障害を測る。
一点で任意に選んだ二本の接方向を、そのまま一般法線に外挿したものではない。
\end{proof}

別に、点 $(a,b,c)=(1,1,1)$ で26個の実族座標を固定した
\emph{全}閉ファイバーを計算する。
その二変数形式表示では非退化な二次式が一本、三次初期式が全三次単項式を張る。
実際の局所環で $\mathfrak m^3/\mathfrak m^4=0$ なので Nakayama の補題により
$\mathfrak m^3=0$ である。
従って全次数で閉ファイバーは
\begin{equation}\label{eq:closed-five}
 \CC[u,v]/(uv,u^3,v^3)
\end{equation}
と同型であり、Hilbert 関数は $(1,2,2)$、長さは5である。
この閉ファイバー計算は全350式と全有理係数で独立に検証されている。

\subsection{有限階数比較から全完備局所環へ}

\begin{proof}[定理\ref{thm:main}(5)の完成環と長さの証明]
まず非零錐上の全126正準係数の完備局所環を $A$ とする。
実際の族の26方向の非零小行列式から、26変数形式環
$R=\CC[[r_1,\ldots,r_{26}]]$ と、合成が恒等である
\[
 R\longrightarrow A\longrightarrow R
\]
を得る。後の矢印は実 Poisson 族への制限である。
式\eqref{eq:closed-five}から完備 Nakayama により $A$ は $R$ 上有限であり、
全射 $R^5\twoheadrightarrow A$ が存在する。
ここではまだ平坦性を仮定しない。

$K=\operatorname{Frac}(R)$ とする。
augmentation が選ぶ一般横断 Artin 局所因子は埋め込み次元2である。
補題\ref{lem:augmented}により二次関係空間の次元は1以下なので、
次数0,1,2の Hilbert 値は少なくとも $(1,2,2)$、長さは5以上となる。
一方、全射 $K^5\twoheadrightarrow A\otimes_RK$ から次元は5以下である。
従って両者は5に等しい。
全射 $R^5\twoheadrightarrow A$ の kernel は自由加群の部分加群で、
一般階数0である。$R$ が整域なので kernel は零、
ここで初めて $A$ が $R$ 自由階数5であると結論できる。

$N=\ker(A\to R)$ とする。
一般横断環の長さ5と Hilbert 値 $(1,2,2)$ から $N_K^3=0$ であり、
$A$ の $R$ 自由性により $N^3=0$ が全環で成立する。
$N/N^2$ は閉ファイバーで二生成、一般階数2なので、同じ自由加群の kernel の議論により
$R$ 自由階数2である。
$N$ は自由階数4で、商 $N/N^2$ が自由なので $N^2$ も自由階数2となる。
この順序で論じることにより、基底変換に伴う Tor を仮定で消していない。

$N/N^2$ の基底を $u,v\in N$ に持ち上げる。
$N^3=0$ により $N^2$ は $u^2,uv,v^2$ で生成され、
$\Sym^2(R^2)\to N^2$ の kernel は自由階数1である。
その二次形式は閉ファイバーで非退化なので、判別式は $R$ の単元となる。
Hensel の平方根と可逆線形変換により関係を $uv$ に移せる。
従って
\[
 R[u,v]/(uv,u^3,v^3)\longrightarrow A
\]
は全射であり、両辺は $R$ 自由階数5なので同型である。
非零錐の自由スカラー因子を一つ除くと、主定理の25滑らか変数の射影表示を得る。

被約環は25変数の形式冪級数環で、極小素は一つである。
Artin 因子の socle は $u^2,v^2$ の張る二次元空間なので、
全完成環は Cohen--Macaulay・type 2 であり、Gorenstein ではない。
一般長さ5と nilradical の三乗零も従う。
\end{proof}

\subsection{パラメータ原点では Cohen--Macaulay でない}

\begin{proof}[定理\ref{thm:main}(6)の証明]
第\ref{sec:upper}節の全八変数横断環を $B$ とすると、
埋め込み次元は8、実次元は3である。
さらに真の五次初期式まで同じ全関係の繰上がりを行うと、
全初期イデアルに含まれる次元3・次数5の斉次イデアル $L_5$ を得る。
この有理係数証明書は独立監査済みであり、四次の次数6から五次の次数5への減少を検出する。
従って Hilbert--Samuel 重複度は $e(B)\le5$ である。

一方、剰余体が無限体である Cohen--Macaulay 局所環には
$e(B)\ge\operatorname{edim}(B)-\dim B+1$ が成立する。
本例なら右辺は $8-3+1=6$ となり矛盾する。
従って $B$ は Cohen--Macaulay でない。
式\eqref{eq:full-formal}の22個の滑らか因子を戻してもこの結論は変わらない。
\end{proof}

点 $[P_{1,1,1}]$ の Cohen--Macaulay 性と $[P_+]$ での失敗は、異なる点での結論である。
後者は全 Poisson 環についての主張であり、被約成分 $\cC$ 自体の大域的正規性や
全特異点を決定したものではない。

\section{追加の幾何と残る問題}

定数二ベクトル $C_0=\partial_y\wedge\partial_z$ は全パラメータで $P$ と両立する。
さらに $a\ne0$ では
\[
 P\left(x,y,z-\frac{3\lambda}{4a},w\right)=P(x,y,z,w)+\lambda C_0
\]
となるので、この pencil は $z$ の平行移動から得られる。
一般パラメータでは定数の両立二ベクトルの空間はこの一次元だけである。
実際、六つの未知定数に対する20本の係数条件は点 $(1,2,3)$ で階数5をもち、
非零五次小行列式は $920192/3$ である。

二次項の線形極化を用いる compact center の条件も線形である。
中心を $[u:v_x:v_y:v_z:v_w]$ と書けば、各括弧の一次部分 $L^{ij}$、二次部分 $Q^{ij}_2$ に対して
\[
 uL^{ij}+\sum_kv_k\partial_kQ^{ij}_2=0
\]
を課す。
全パラメータで $[0:0:0:1:0]$ が解であり、一般には唯一の解である。
証人となる24行5列行列は点 $(1,2,3)$ で階数4、小行列式は $-1245184/3$ である。
この中心は無限遠にある。
有限中心の存在を仮定した分類に置き換えることはできない。

本稿では実際の成分と、その成分に沿う全スキームの一般非被約性を証明した。
一般横断環は長さ5と決定され、指定一般点での全乗法も上記の完成環表示で与えられた。
また $P_+$ における全局所次元25の証明は、全形式枝の本数や完全初期イデアルを決定しない。
今後の課題には、二変数表示に残る有限同値の決定、一次像超曲面の方程式、
本成分の正規性と特異点、他成分との交叉が含まれる。
射影 Poisson 構造の一般的背景については\cite{pym}を参照されたい。

\appendix
\section{係数順序と26次小行列式}\label{app:rank}

添字は0から数える。行順序は第4節の150係数の順序とする。
列順序は辞書式の $E_{ij}$ 25本、スカラー、$\partial_a,\partial_b,\partial_c$ とする。
補題\ref{lem:rank}の26次小行列式で用いた行は
\[
\begin{gathered}
0,1,2,4,15,16,17,19,30,31,32,34,45,46,60,61,64,\\
75,76,79,90,105,106,109,110,113
\end{gathered}
\]
であり、列は
\[
0,1,2,3,4,5,6,7,8,9,10,11,12,13,14,15,16,18,19,20,21,22,23,25,26,27
\]
である。末尾の族方向 $\partial_c$ を使わずとも非零である。
この添字は保存された全列従属関係と合わせて、階数下界の小さな有理証人になる。

\section{計算証明書と再現手順}\label{app:cert}

証明に用いた等式は全て有理数体または表示した有理関数体上の厳密等式である。
有限体の標本のみを根拠とする次元上界は本稿では使っていない。
計算資源の都合で方程式を捨てる操作と、可逆行列・形式陰関数定理による同値な消去を区別する。

ファイル名の共通末尾を $\tau=\texttt{20261103\_v63}$ と略記する。
以下は研究リポジトリの相対パスである。
\begin{longtable}{>{\raggedright\arraybackslash}p{0.23\textwidth}>{\raggedright\arraybackslash}p{0.70\textwidth}}
\toprule
内容 & ファイルの幹（コードは \texttt{code/}、ノートは \texttt{research/notes/}）\\
\midrule\endhead
明示族・全 Jacobi・階数26
 & \texttt{q\_positive\_weight6\_new\_family\_}\(\tau\)、
   独立検証 \texttt{q\_positive\_weight6\_new\_family\_review\_}\(\tau\)\\
全八変数形式模型
 & \texttt{q\_positive\_weight6\_quadratic\_model\_}\(\tau\)\\
真の四次初期式と次元3
 & \texttt{q\_positive\_weight6\_initial4\_}\(\tau\)、
   独立検証 \texttt{q\_positive\_weight6\_small\_model\_review\_}\(\tau\)\\
真の五次初期式・重複度上界5
 & \texttt{q\_positive\_weight6\_initial5\_}\(\tau\)、
   \texttt{Q\_POSITIVE\_WEIGHT6\_INITIAL5\_}\(\tau\) の保存結果と独立検証\\
成分25の認識
 & \texttt{q\_positive\_weight6\_component25\_review\_}\(\tau\)\\
実因子・モジュラー流・座標変換
 & \texttt{r4\_positive\_weight6\_geometry\_}\(\tau\)、
   全斉次変換の別検査 \texttt{r4\_positive\_weight6\_transform\_audit\_}\(\tau\)\\
toric 実因子による相異性
 & \texttt{r4\_positive\_weight6\_distinctness\_}\(\tau\)\\
一次ジェット関数体階数
 & \texttt{q\_new\_family\_firstjet\_image\_}\(\tau\)、
   独立検証 \texttt{q\_new\_family\_firstjet\_review\_}\(\tau\)\\
全 Jacobi 関数体階数98
 & \texttt{q\_new\_family\_generic\_jacobian\_rank\_}\(\tau\)、
   独立検証 \texttt{q\_new\_family\_generic\_jacobian\_review\_}\(\tau\)\\
二法線二次障害・全閉 Artin fiber
 & \texttt{q\_new\_family\_generic\_quadratic\_rank\_}\(\tau\)、
   独立検証 \texttt{q\_new\_family\_generic\_quadratic\_review\_}\(\tau\)、
   \texttt{q\_new\_family\_transverse\_artin\_review\_}\(\tau\)\\
全完成環・正点の非CM性
 & \texttt{q\_new\_family\_length5\_completion\_}\(\tau\)、
   \texttt{q\_positive\_weight6\_noncm\_}\(\tau\)\\
\bottomrule
\end{longtable}

主要な実行結果は対応する大文字の幹を持つ \texttt{reports/} 下にあり、
入力の SHA--256、選択小行列、全従属関係を固定している。
特に次の圧縮 JSON が、全方程式から初期式に至る有限証明書である。
\begin{quote}\small\ttfamily
reports/Q\_POSITIVE\_WEIGHT6\_QUADRATIC\_MODEL\_20261103\_v63/certificate.json.gz\\
reports/Q\_POSITIVE\_WEIGHT6\_INITIAL4\_20261103\_v63/certificate.json.gz
\end{quote}
後者の SHA--256 は
\begin{quote}\footnotesize\ttfamily
0c30e393e7de146949f42ba2acb91f49d6c1bf88fe2311e57ecf46c94add8a38
\end{quote}
である。
五次初期式の全繰上がりは
\path{reports/Q_POSITIVE_WEIGHT6_INITIAL5_20261103_v63/fifth-complete-spanning.json.gz}
に保存し、独立検証は対応する \texttt{\_REVIEW/report.json} に記録する。
21個の初期生成元と Hilbert 級数の検証用入力は、この論文専用の補助ファイル
\texttt{positive\_weight6\_initial21.m2} にも書き出す。
四次項の長い展開を省略しても、元の Jacobi 方程式に戻す係数証明書は保持されている。

代表的な再現コマンドは次のとおりである。
既存の Cloud 環境の実行環境を使い、出力先を新規ディレクトリに指定する。
\begin{quote}\begin{minipage}{\linewidth}\footnotesize\ttfamily
source /workspace/cloud-setup/activate.sh\\
.venv/bin/python code/r4\_positive\_weight6\_geometry\_20261103\_v63.py
 --output /tmp/positive-geometry-replay\\
.venv/bin/python code/q\_positive\_weight6\_initial4\_20261103\_v63.py
 --output /tmp/positive-initial4-replay\\
.venv/bin/python code/q\_new\_family\_firstjet\_review\_20261103\_v63.py
 --output /tmp/positive-firstjet-review
\end{minipage}\end{quote}
独立検証器のコマンド引数や依存先は、各ファイルの \texttt{--help} と対応ノートに記録する。
原稿・補助データ・PDF のハッシュと実測ページ数は
\path{reports/POSITIVE_WEIGHT6_COMPONENT_PAPER_20261103_v63/}
に保存する。

\begin{thebibliography}{9}
\bibitem{mps}
M.~Matviichuk, B.~Pym, T.~Schedler,
\emph{A local Torelli theorem for log symplectic manifolds},
arXiv:2010.08692v2 (2021), 特に Theorem~1.4, Theorem~6.16, Example~6.24.
\bibitem{pym}
B.~Pym,
\emph{Constructions and classifications of projective Poisson varieties},
Letters in Mathematical Physics \textbf{108} (2018), 573--632;
arXiv:1701.08852.
\bibitem{repo}
Poisson P4 研究プロジェクト,
\emph{研究ノートと厳密計算証明書}, v63 (2026年11月3日),
\url{https://github.com/uedakazushi/poisson-p4}.
本稿の初期式・有理関数体階数・旧成分カタログの具体的入力は付録\ref{app:cert}を参照。
\end{thebibliography}
\end{document}
